% WACV 2027 -- PhysWeep supplementary material.
% Standalone document: compiles on its own into a separate PDF. Shares
% preamble.tex with main.tex so notation matches exactly.
%
% Section numbering here is S1, S2, ... because main.tex refers to
% "Sec. S1/S2/S3" by name (the two documents cannot share \label/\ref).
% If you renumber or reorder anything here, update those pointers in main.tex.

\documentclass[10pt,twocolumn,letterpaper]{article}

% \usepackage[review,algorithms]{wacv}
% \usepackage[review,applications]{wacv}
\usepackage[review,datasets]{wacv}
% \usepackage{wacv}

% ============================================================================
% PhysWeep preamble (shared by main.tex and appendix.tex).
% wacv.sty already provides: times, xcolor (dvipsnames), graphicx, amsmath,
% amssymb, booktabs, natbib[numbers,sort&compress], caption, subcaption, url,
% enumitem[shortlabels,inline]. Do NOT reload any of those here.
% ============================================================================

\usepackage{amsthm}
\usepackage{array}
\usepackage{tikz}
\usetikzlibrary{arrows.meta,positioning,calc,backgrounds,fit,patterns}
\usepackage{algorithm}
\usepackage{algpseudocode}
\usepackage{pifont}

% --- red-flag macros: every item the authors must supply before submission.
% Search the source for "MISSING" / "VERIFY" to find them all; they render in
% red in the PDF so none can ship unnoticed.
\newcommand{\MISSING}[1]{\textcolor{red}{[MISSING: #1]}}
\newcommand{\VERIFYFLAG}[1]{\textcolor{red}{[VERIFY: #1]}}

\newcommand{\cmark}{\textcolor{green!55!black}{\ding{51}}}
\newcommand{\xmark}{\textcolor{red!70!black}{\ding{55}}}

% Keep theorem statements from orphaning their header across a column break.
% Without this, a Proposition whose header lands exactly at a column boundary
% is shipped out above the running header and bleeds off the page edge
% (observed as "Overfull \hbox ... while \output is active").
\usepackage{needspace}
\usepackage{colortbl}   % \cellcolor for the coverage heat map
\usepackage{multirow}

% --- coverage heat map -----------------------------------------------------
% Single-hue sequential ramp so the table stays legible in grayscale, which
% top venues still print in. Zero is rendered in gray rather than white so a
% measured zero is visually distinct from an empty cell.
\definecolor{trkhi}{RGB}{31,102,160}
\newcommand{\tzero}{\cellcolor{black!8}\textcolor{black!55}{$0$}}
\newcommand{\tk}[2]{\cellcolor{trkhi!#1}$#2$}
\newcommand{\tkb}[2]{\cellcolor{trkhi!#1}$\mathbf{#2}$}
% compact in-table placeholder for a value the authors must still compute;
% the matching \MISSING{} lives in the caption so it reaches the red-flag list
\newcommand{\tbd}{\cellcolor{red!12}\textcolor{red!80!black}{\textbf{?}}}

% Non-floating tables in a plain center environment inherit \topsep and
% \parskip above and below, which costs roughly 25pt per table. This wrapper
% keeps the centring and drops the padding.
\newenvironment{tightcenter}
  {\setlength{\topsep}{1pt}\setlength{\partopsep}{0pt}%
   \setlength{\parskip}{0pt}\par\centering}
  {\par\addvspace{2pt}}
\newtheorem{proposition}{Proposition}
% Unnumbered form for restating a main-paper result in the supplementary, so
% the supplementary does not assign it a second, different number.
\newtheorem*{propositionstar}{Proposition}
\newtheorem{definition}{Definition}
\newtheorem{remark}{Remark}
\newtheorem{assumption}{Assumption}

\newcommand{\PRE}{\ensuremath{\mathrm{PRE}}}
\newcommand{\PRI}{\ensuremath{\mathrm{PRI}}}
\newcommand{\CBI}{\ensuremath{\mathrm{CBI}}}
\newcommand{\EG}{\ensuremath{\mathrm{EG}}}
\newcommand{\thetah}{\ensuremath{\hat{\theta}}}

% --- line-number overlap fix (applies only if wacv.sty loaded lineno)
\makeatletter
\AtBeginDocument{\@ifpackageloaded{lineno}{%
  \setlength{\linenumbersep}{6pt}%
  \renewcommand\linenumberfont{\normalfont\scriptsize\color{gray}}%
  % Consecutive line numbers can stack on top of each other where a
  % display or a centered block compresses baseline spacing. Setting a
  % minimum inter-number distance stops that without touching text layout.
  \setlength{\linenumberwidth}{18pt}%
  \leftlinenumbers}{}}
\makeatother


\definecolor{wacvblue}{rgb}{0.21,0.49,0.74}
\usepackage[pagebackref,breaklinks,colorlinks,allcolors=wacvblue]{hyperref}

\def\wacvPaperID{2935} % must match main.tex
\def\confName{WACV}
\def\confYear{2027}

% Must match main.tex.
\title{\textbf{\textsc{PhysWeep}: Does a Video Generator Realize\\[2pt]
the Physics You Ask For?}}
\author{Anonymous Authors\\ Under double-blind review at WACV 2027}

\renewcommand{\thesection}{S\arabic{section}}

\begin{document}
\maketitlesupplementary

% ---------------------------------------------------------------------------
\section{Measurement pipeline and failure-form estimation}
\label{app:pipeline}
% ---------------------------------------------------------------------------

Algorithm~\ref{alg:pipeline} gives the full measurement loop referenced from
the main paper's Suite, Pipeline, and Protocol section.

\begin{algorithm}[H]
\caption{\textsc{PhysWeep} measurement, one system and one model}
\label{alg:pipeline}
\begin{algorithmic}[1]
\Require frozen generator $G$; sweep grid $\{\theta_i\}$; seeds $\{z_j\}$;
         law $f$; tracker $T$; fit-quality gate $\rho$; VLM rater $P$
\For{each $\theta_i$ and seed $z_j$}
  \State $c_{ij}\gets \textsc{Render}(s_0(z_j),\theta_i)$
  \State $v_{ij}\gets G(c_{ij})$ \Comment{inference only, weights frozen}
  \State $o_{ij}\gets T(v_{ij})$ \Comment{blob centroid or CoTracker3}
  \State $(\thetah_{ij}, R^2_{ij})\gets \textsc{Fit}(f, o_{ij})$
  \If{$R^2_{ij}<\rho$}
    \State mark $v_{ij}$ \textbf{non-physical}; exclude from \PRE
  \EndIf
  \State $p_{ij}\gets P(v_{ij})$ \Comment{plausibility contrast}
\EndFor
\State compute \PRE, slope $\beta$, Spearman $\rho_s$, \EG, \PRI, \CBI
\State bootstrap $95\%$ CIs over $(\theta_i,\text{seed})$ pairs
\State \Return metrics, gate-dropped rate, plausibility pairs
\end{algorithmic}
\end{algorithm}

\paragraph{Prior reversion (shrinkage, \PRI).}
Given pairs $(\theta_i,\thetah_{ij})$, fit
$\thetah=\alpha\theta+(1-\alpha)\theta_0$ by least squares with $\theta_0$
either fixed to the physically typical value or estimated jointly;
$\PRI=1-\hat\alpha$, clipped to $[0,1]$.

\paragraph{Case clamping (\CBI).}
Fit $\thetah=\min(\theta,\theta_{\text{edge}})$ with $\theta_{\text{edge}}$
at the in-range boundary or estimated;
$\CBI=1-\mathrm{RSS}_{\text{clamp}}/\mathrm{RSS}_{\text{null}}$.

\paragraph{Selection.}
Compare prior reversion and case clamping by leave-one-$\theta$-out $R^2$ and by $\Delta$AIC. A
hypothesis is selected only if its held-out $R^2$ is positive; if both fail
that test the pair is reported as \emph{neither} and the observed behavior is
described without forcing a label. This is what happens on all eight
adjudicated pairs in the main paper, for the reason given there. All
confidence intervals use the paired bootstrap over $(\theta_i,\text{seed})$.

\paragraph{When \CBI\ is not identifiable.}
The released smoke test records a limitation of \CBI\ as defined. If an
out-of-range grid lies entirely on one side of the in-range band, the clamp
prediction $\min(\theta,\theta_{\text{edge}})$ is constant across
$S_\text{out}$ and is therefore indistinguishable from the null mean model,
so $\CBI\approx0$ and clamping cannot be adjudicated at any sample size.
Pendulum damping is the axis where this specific degeneracy applies: its
out-of-range grid $\{0.3,0.6,0.9\}$ lies entirely above the in-range grid;
no result in the main paper depends on this, because every pendulum row is
gate-dropped at $96.7$--$100\%$ and carries no adjudication. A two-sided
damping grid would be needed to test clamping on that axis.
Separately, \texttt{select\_mechanism}'s default edge
($\theta_{\text{edge}}=\min(S_\text{out})$) is a crude choice even on our
two-sided axes: on data whose recovered values are seed-clustered rather
than clamped, this default can still fit only slightly worse than the null
mean and clip to $\CBI\approx0$ rather than surface the (often more
negative) fit quality a properly chosen edge would show. This does not
change any adjudication in the main paper, since clamping is rejected on
every pair either way (its held-out $R^2$ is non-positive throughout, on
all eight adjudicated pairs in the main paper's mechanism-adjudication
table), but it means a near-zero \CBI\ on our data should be read as
``clamping does not fit,'' not as evidence the edge choice was
well-identified.

% ---------------------------------------------------------------------------
\section{Consistency of the recovery estimator: full statement and proof}
\label{app:proof}
% ---------------------------------------------------------------------------
We restate the main paper's consistency proposition with explicit regularity
conditions and give the argument in full. The claim concerns only the
estimator $R$, not the generator $G$.

\paragraph{Setup.}
Fix a system with law $f$ and scalar parameter $\theta\in\Theta$, with
$\Theta$ a compact interval. A clip of $T$ frames at a given resolution
yields a tracked observable $o\in\mathbb{R}^{d}$ (object centers, angles, or
contact times). Write the population fitting statistic as
$\mathcal{T}(\theta)\in\mathbb{R}^{k}$: for the projectile, the quadratic
coefficient of height against time; for the pendulum, the mean zero-crossing
interval; for the bouncing ball, the successive height ratio. The estimator
forms an empirical statistic $\hat{\mathcal{T}}$ from $o$ and returns
$\thetah=g(\hat{\mathcal{T}})$, where $g=\mathcal{T}^{-1}$ on the range of
$\mathcal{T}$.

\paragraph{Assumptions.}
\textbf{(Trackability)} the tracker returns
$\hat o = o^\star + \varepsilon$ with $\mathbb{E}[\varepsilon]=0$ and
$\operatorname{Var}(\varepsilon)=\sigma^2\Sigma$ for a fixed positive-definite
$\Sigma$, and $\sigma^2\to 0$ as $T$ and the resolution grow.
\textbf{(Identifiability)} $\mathcal{T}:\Theta\to\mathcal{T}(\Theta)$ is
injective and continuously differentiable with
$\|\mathcal{T}'(\theta)\|\ge c>0$ on $\Theta$.
\textbf{(Regularity)} the empirical statistic is a continuous functional
$\hat{\mathcal{T}}=\Phi(\hat o)$ with $\Phi$ continuous at $o^\star$ and
$\Phi(o^\star)=\mathcal{T}(\theta)$.

\begin{propositionstar}[Consistency of $R$, restating the main paper]
\label{prop:consistency-restated}
Under trackability, identifiability, and regularity,
$\thetah=g(\Phi(\hat o)) \xrightarrow{p} \theta$ as $\sigma^2\to 0$.
\end{propositionstar}

\begin{proof}
By identifiability, $\mathcal{T}$ is a continuous injection on the compact set $\Theta$,
so its inverse $g=\mathcal{T}^{-1}$ is continuous on $\mathcal{T}(\Theta)$: a
continuous bijection from a compact space to a Hausdorff space is a
homeomorphism. By trackability, $\operatorname{Var}(\hat o)=\sigma^2\Sigma\to 0$ with
mean $o^\star$, so $\hat o\to o^\star$ in mean square and hence in
probability. By regularity, $\Phi$ is continuous at $o^\star$, so the continuous
mapping theorem gives
$\hat{\mathcal{T}}=\Phi(\hat o)\xrightarrow{p}\Phi(o^\star)=\mathcal{T}(\theta)$.
Applying the continuous mapping theorem again with the continuous $g$ yields
$\thetah=g(\hat{\mathcal{T}})\xrightarrow{p} g(\mathcal{T}(\theta))=\theta$.
\end{proof}

\paragraph{Per-system check of identifiability and regularity.}
\emph{Projectile.} $\mathcal{T}(\theta)=-\theta/2$ is linear, hence injective
with $\mathcal{T}'=-1/2$, and $\Phi$ is ordinary least squares on a quadratic
basis, continuous in the data.

\emph{Pendulum, and a bias this exposes.} For the undamped case $\zeta=0$,
$\mathcal{T}(\omega)=2\pi/\omega$ on $\omega>0$ is strictly monotone with
nonvanishing derivative on any $\Theta$ bounded away from $0$, and
zero-crossing interpolation is continuous where crossings are simple. Our
pendulum is damped ($\zeta>0$ throughout,
Section~\ref{app:systems}), so the statistic realized by zero-crossing
spacing is $\mathcal{T}(\omega,\zeta)=2\pi/(\omega\sqrt{1-\zeta^2})$, the
\emph{damped} frequency $\omega_d$, not the natural frequency $\omega$. The
map remains injective in $\omega$ for fixed $\zeta$, so
Proposition~\ref{prop:consistency-restated} still guarantees consistent
recovery of $\omega_d$; the reported $\hat\omega$ is therefore a slightly
biased estimate of $\omega$ whenever $\zeta>0$. Verified numerically, the
bias is $<1.3\%$ at the in-range $\zeta$ grid $\{0.02,0.05,0.1\}$, grows with
$\zeta$, and is materially larger at the out-of-range grid
$\{0.3,0.6,0.9\}$. We did not correct for it post hoc, since doing so would
change every already-reported $\hat\omega$; we state it so the recovery is
not mistaken for exact, and note it compounds with the heavily-damped,
few-peaks low-fidelity regime for this axis. In practice this caveat has
little effect on the reported results, because every pendulum row in the main
paper is gate-dropped at $96.7$--$100\%$ and carries no slope estimate.

\emph{Bouncing ball.} The map $e\mapsto e^2$ from restitution to successive
height ratio is injective and smooth on $(0,1]$.

The fit-quality gate ($R^2<\rho$) excludes generations on which $\Phi$ is
ill-posed, for example when the object morphs or vanishes, so regularity
is enforced operationally rather than assumed.

% ---------------------------------------------------------------------------
\section{Seed-conditional defaults: full statement and proof}
\label{app:seed}

This section proves the main paper's Proposition on misspecification and
records the estimator for the Seed Determination Index.

\textbf{Model.} The sweep is a \emph{crossed} design: every grid value
$\theta_i$ is run with every seed $z_j$, so the seed is independent of
$\theta$ by construction. Under a seed-conditional default,
$\thetah_{ij}=c(z_j)+\varepsilon_{ij}$ with $\mathbb{E}[\varepsilon]=0$,
$\operatorname{Var}(\varepsilon)=\sigma_\varepsilon^2$,
$\varepsilon\perp(\theta,z)$, and
$\sigma_c^2=\operatorname{Var}_z[c(z)]>0$. The crossing assumption is load
bearing and we flag it for anyone reusing the protocol. If seeds were
confounded with grid position, so that $c$ varied systematically with
$\theta$, statement (i) fails and the failure is not subtle: in simulation a
generator whose output is entirely seed-determined, with no response to
$\theta$ whatsoever, returns $\hat\alpha=1.006$ and $\PRI=-0.006$ under a
confounded design, which reads as \emph{perfect faithfulness}. Crossing seeds
with grid values is therefore a correctness requirement of the measurement,
not a convenience.

\begin{propositionstar}[Restating Proposition~1 of the main paper]
Under this model: (i) the pooled OLS fit of $\thetah=\alpha\theta+b$
satisfies $\hat\alpha\xrightarrow{p}0$, so $\PRI\to1$; (ii)
$\hat b\xrightarrow{p}\mathbb{E}_z[c(z)]$ and the implied
$\hat\theta_0=\hat b/(1-\hat\alpha)\xrightarrow{p}\mathbb{E}_z[c(z)]$; (iii)
any measurable $h(\theta)$ has population $R^2=0$ for predicting $\thetah$,
its leave-one-$\theta$-out estimate satisfies
$R^2_{\mathrm{LOO}}\xrightarrow{p}0$, and for any fixed $\tau>0$ a rule
accepting only $R^2_{\mathrm{LOO}}>\tau$ returns \emph{neither} with
probability tending to one. None of (i)--(iii) uses $\sigma_c^2>0$, so all
three hold verbatim when $\sigma_c^2=0$ and $\sigma_\varepsilon^2$ is large.
\end{propositionstar}

\begin{proof}
(i) $\operatorname{Cov}(\theta,\thetah)=\operatorname{Cov}(\theta,c(z))
+\operatorname{Cov}(\theta,\varepsilon)=0$: the first term vanishes because
the crossed design makes $z$ independent of $\theta$, the second because
$\varepsilon\perp\theta$. The OLS slope is
$\widehat{\operatorname{Cov}}(\theta,\thetah)/\widehat{\operatorname{Var}}(\theta)$,
which converges in probability to $0/\operatorname{Var}(\theta)=0$ by the
weak law, with $\operatorname{Var}(\theta)>0$ since the grid is
non-degenerate. Hence $\PRI=1-\hat\alpha\to1$.

(ii) $\hat b=\overline{\thetah}-\hat\alpha\bar\theta\xrightarrow{p}
\mathbb{E}[\thetah]=\mathbb{E}_z[c(z)]$, and dividing by $1-\hat\alpha\to1$
gives the same limit. This limit is a population average: when $c$ takes two
well-separated values it need not lie near either, so the fitted ``default''
describes no individual generation.

(iii) For any measurable $h$,
$\mathbb{E}[(\thetah-h(\theta))^2]
=\operatorname{Var}(\thetah)
+\mathbb{E}[(\mathbb{E}[\thetah]-h(\theta))^2]$,
the cross term vanishing by $\thetah\perp\theta$. This is minimized at
$h\equiv\mathbb{E}[\thetah]$ with value $\operatorname{Var}(\thetah)$, so the
population $R^2$ is $0$. The leave-one-$\theta$-out numerator and denominator
are averages of i.i.d.\ terms converging in probability to
$\operatorname{Var}(\thetah)>0$, so their ratio converges to $1$ and
$R^2_{\mathrm{LOO}}\xrightarrow{p}0$ by the continuous mapping theorem.
Consequently $\Pr[R^2_{\mathrm{LOO}}>\tau]\to0$ for any fixed $\tau>0$, which
is the claim. Finally, no step above invokes $\sigma_c^2>0$; the argument
uses only $\thetah\perp\theta$, which also holds when $c$ is constant and the
variance is carried entirely by $\varepsilon$.
\end{proof}

\begin{remark}[What is deliberately not claimed]
We do \emph{not} claim $\mathbb{E}[R^2_{\mathrm{LOO}}]<0$. Simulation under
this model gives a mean of $-0.024$ at a $5\times5$ design shrinking to
$-0.001$ at $20\times20$, so the expectation does appear non-positive and to
vanish with sample size, but $R^2_{\mathrm{LOO}}$ is a ratio of dependent
random variables and we have not discharged the argument rigorously. Nothing
in the paper depends on it: convergence in probability to $0$, together with
the fixed margin, is what statement (iii) and the released selection rule
actually use. We record it here as an observation rather than a result.
\end{remark}

\begin{remark}[Why a margin is required in the selection rule]
Because $R^2_{\text{LOO}}$ is centred on zero under this model, its
\emph{sign} is noise. An acceptance test of the form $R^2_{\text{LOO}}>0$
therefore labels unstructured data as prior reversion with probability
approaching $1/2$. The released \texttt{select\_mechanism} requires
$R^2_{\text{LOO}}>0.1$; the smoke test contains a synthetic sample on which
the bare sign test returns \texttt{prior\_reversion} ($R^2_{\text{LOO}}=0.035$)
while the margin rule correctly returns \texttt{neither}, and asserts both, so
the margin cannot be removed silently.
\end{remark}

\paragraph{Estimating \texorpdfstring{$\mathrm{SDI}$}{SDI}.}
With $k$ seeds, $N$ gated clips and mean group size $n_0=N/k$, let
$\mathrm{MS}_b$ and $\mathrm{MS}_w$ be the between- and within-seed mean
squares of $\thetah$. We report the one-way random-effects estimate
\begin{equation*}
\mathrm{SDI}=\frac{\mathrm{MS}_b-\mathrm{MS}_w}
                  {\mathrm{MS}_b+(n_0-1)\,\mathrm{MS}_w},
\end{equation*}
clipped to $[0,1]$, which is consistent for
$\sigma_c^2/(\sigma_c^2+\sigma_\varepsilon^2)$. The uncorrected variance
share is strongly upward-biased at this design size, returning about $0.18$
in simulation where the truth is $0$, which is precisely the regime in which
a false positive would be most damaging; the corrected estimate returns
about $0.07$ there. $\mathrm{FDI}$ is the analogous share for the grid factor.
Both return \texttt{NaN}, never a number, when fewer than two seeds are
present, when $\thetah$ is constant, or when there is no within-seed
replication.

\section{System grids, seed counts, and rendering}
\label{app:systems}
% ---------------------------------------------------------------------------
All scenes are rendered by a seeded deterministic 2D engine at
$256\times256$ over $T{=}24$ frames at $24$\,fps, with a single high-contrast
disk (radius $8$\,px) on a plain background and a fixed ground line. Only the
swept parameter and the seeded initial condition vary within a system.
Spatial and temporal scales are held fixed within a clip and all metrics use
the ratio $\thetah/\theta$, so the result is invariant to a common
rescaling. The in-range split is centered on values common in natural video;
the out-of-range split probes rare or absent values. These are the grids
fixed in the released config and used for every reported result.

\begin{center}
\renewcommand{\arraystretch}{1.05}
\resizebox{\columnwidth}{!}{%
\begin{tabular}{@{}llll@{}}
\toprule
\textbf{System} & \textbf{$\theta$} & \textbf{In-range grid} & \textbf{Out-of-range grid} \\
\midrule
Projectile      & $g$        & $\{5,7,9.8,12,15\}$         & $\{1.6,\,3,\,20,\,25\}$ \\
Damped pendulum & $\omega$   & $\{2,3,4,5\}$ rad/s         & $\{0.5,\,1,\,8,\,12\}$ \\
                & $\zeta$    & $\{0.02,0.05,0.1\}$         & $\{0.3,\,0.6,\,0.9\}$ \\
Bouncing ball   & $e$        & $\{0.6,0.7,0.8,0.9\}$       & $\{0.2,\,0.35,\,0.97\}$ \\
Spring--mass    & $k$        & $\{20,40,60,80\}$ N/m       & $\{5,\,10,\,200,\,400\}$ \\
Inclined slide  & $\mu$      & $\{0.1,0.2,0.3,0.4\}$       & $\{0.01,\,0.7,\,1.0\}$ \\
\bottomrule
\end{tabular}}
\captionof{table}{Sweep grids per system. Each row has one swept axis;
$m\ge5$ seeded initial conditions per grid point.}
\label{tab:grids}
\end{center}

\paragraph{Seed counts.}
The main paper refers to this section for the exact seed counts behind each
reported cell, including the targeted follow-up batch run on the two
model-axis pairs closest to \texttt{select\_mechanism}'s four-point
leave-one-out threshold.

\begin{center}
\renewcommand{\arraystretch}{1.05}
\resizebox{\columnwidth}{!}{%
\begin{tabular}{@{}llccc@{}}
\toprule
\textbf{Model} & \textbf{Axis} & \textbf{seeds/pt (in/out)} & \textbf{clips gen.} & \textbf{clips gated} \\
\midrule
LTX-Video & projectile & $5/10$ & $55$ & $12$ \\
LTX-Video & pendulum $\omega$ & $5/55$ & $130$ & $1$ \\
LTX-Video & pendulum $\zeta$ & $5/15$ & $60$ & $0$ \\
LTX-Video & bouncing ball & $5/15$ & $65$ & $5$ \\
LTX-Video & spring-mass & $5/15$ & $80$ & $0$ \\
LTX-Video & inclined slide & $5/5$ & $35$ & $11$ \\
\addlinespace[2pt]
CogVideoX & projectile & $5/5$ & $40$ & $32$ \\
CogVideoX & pendulum $\omega$ & $5/13$ & $46$ & $0$ \\
CogVideoX & pendulum $\zeta$ & $5/13$ & $54$ & $0$ \\
CogVideoX & bouncing ball & $5/33$ & $119$ & $10$ \\
CogVideoX & spring-mass & $5/13$ & $72$ & $0$ \\
CogVideoX & inclined slide & $5/13$ & $59$ & $0$ \\
\addlinespace[2pt]
DynamiCrafter & projectile & $5/10$ & $55$ & $33$ \\
DynamiCrafter & pendulum $\omega$ & $5/5$ & $30$ & $0$ \\
DynamiCrafter & pendulum $\zeta$ & $5/5$ & $30$ & $1$ \\
DynamiCrafter & bouncing ball & $5/10$ & $50$ & $11$ \\
DynamiCrafter & spring-mass & $5/5$ & $40$ & $0$ \\
DynamiCrafter & inclined slide & $5/10$ & $50$ & $36$ \\
\bottomrule
\end{tabular}}
\captionof{table}{Seed counts and clip yields per model-axis pair. ``seeds/pt
(in/out)'' is the number of seeds run at each in-range and out-of-range grid
value respectively; most pairs use the base $5$ seeds on both splits, but
pairs close to \texttt{select\_mechanism}'s four-gated-point leave-one-out
threshold received a larger targeted follow-up batch on the out-of-range
split only (LTX-Video/pendulum-$\omega$: extended to $55$ seeds/point,
i.e.\ $110$ out-of-range attempts across its two out-of-range grid values,
still yielding only $1$ gated clip; CogVideoX/bouncing-ball: extended to
$33$ seeds/point). No in-range seed count was extended beyond the base
$5$, since the out-of-range split is what \texttt{select\_mechanism}'s
adjudication gate requires.}
\end{center}

\paragraph{Models, prompts, and compute.}
\begin{table*}[t]
\centering
\small
\setlength{\tabcolsep}{5pt}
\renewcommand{\arraystretch}{1.1}
\resizebox{\textwidth}{!}{%
\begin{tabular}{@{}llll@{}}
\toprule
\textbf{Model} & \textbf{Repository} & \textbf{Licence} & \textbf{Loaded via} \\
\midrule
LTX-Video & \texttt{Lightricks/LTX-Video} & LTX-Video Open Weights (custom) & \texttt{diffusers==0.39.0} \\
CogVideoX-5B-I2V & \texttt{zai-org/CogVideoX-5b-I2V} & CogVideoX custom licence\textsuperscript{a} & \texttt{diffusers==0.39.0} \\
DynamiCrafter\textsubscript{512} & \texttt{Doubiiu/DynamiCrafter\_512} & weights: research/non-commercial\textsuperscript{b} & standalone repo, subprocess \\
Qwen2.5-VL-7B-Instruct & \texttt{Qwen/Qwen2.5-VL-7B-Instruct} & Apache 2.0 & \texttt{transformers==4.57.6} \\
CoTracker3 & \texttt{facebookresearch/co-tracker} & CC-BY-NC 4.0\textsuperscript{c} & \texttt{torch.hub.load} \\
\bottomrule
\end{tabular}}
\caption{Every model used, all frozen and inference-only, loaded with no
explicit commit/revision pin -- each resolved to that repository's
default-branch head at run time, a gap against bit-for-bit
reproducibility we disclose rather than paper over.
\textsuperscript{a}Free for academic/research use; commercial use needs
separate registration. \textsuperscript{b}Weights are
research/non-commercial; the companion code repo is separately
Apache-2.0 (inference here uses the weights). \textsuperscript{c}Non-commercial;
some unused sub-components carry MIT or Apache 2.0 instead.}
\end{table*}

\paragraph{Plausibility prompt (verbatim, identical for all $630$ ratings).}
\begin{quote}\small
``You are watching a short video clip of a physics simulation (a ball or
similar object moving under some physical process such as falling,
swinging, bouncing, or sliding). Rate how PHYSICALLY PLAUSIBLE the motion
looks, on a scale from 1 to 5, where 1 means the motion looks clearly
unnatural, broken, or impossible (e.g. teleporting, morphing, floating
with no cause), and 5 means the motion looks completely natural and
physically believable, like something you could film in the real world.
Respond with ONLY a single digit from 1 to 5, nothing else.''
\end{quote}

\paragraph{Prompt-phrasing ablation templates.} Both variants are
per-system, per-$\theta$ templates (not a single fixed string), evaluated
against the same default template used for every other result. The
\emph{terse} variant states only the parameter name and numeric value with
no qualitative anchor phrase, e.g. for projectile: ``a dark ball,
acceleration \{$\theta$:.2f\} meters per second squared''. The
\emph{verbose} variant appends redundant physical-accuracy framing to the
default prompt: ``\{default prompt\}, precise physically accurate
simulation, exact numerical value, photorealistic rendering, laboratory
measurement setup''.

\paragraph{Classifier-free guidance values swept.} $1.0$, $1.5$, $4.5$,
$9.0$, on the projectile system, for LTX-Video and CogVideoX; skipped for
DynamiCrafter, whose guidance scale is hardcoded inside its own inference
subprocess call and not exposed as a pass-through parameter.

\paragraph{Compute.} All generation and analysis ran on NVIDIA A100-SXM4
GPUs (40\,GB) on a shared cluster, inference-only throughout (no training
or fine-tuning at any point). Peak VRAM at load: LTX-Video $14.25$\,GB
(fits a 24\,GB consumer card); CogVideoX-5B-I2V $25.3$\,GB with VAE tiling
enabled (does not). Approximate wall-clock generation time per clip:
LTX-Video $\sim\!2$--$6$\,s; CogVideoX $\sim\!2.3$--$5.9$\,min, varying
with clip length across systems; DynamiCrafter $\sim\!70$--$85$\,s (a
fresh subprocess per clip, so this is dominated by per-call checkpoint
load, not sampling time alone).

\textbf{Total wall-clock GPU-hours: $\sim\!57$ (estimated, not a full
audit).} $54.4$\,h is a direct sum over the $69$ of $\sim\!134$ result
files that log their own \texttt{elapsed\_seconds} per run, covering the
large majority of clips generated across the project. The remaining
$\sim\!95$ clips (the CogVideoX and DynamiCrafter projectile pilots,
which predate this project's practice of logging per-job timing) are
estimated at $\sim\!2.6$\,h by applying each model's own measured
mean seconds-per-clip (from the $69$ logged files: LTX $3.96$\,s/clip,
CogVideoX $122.45$\,s/clip, DynamiCrafter $79.19$\,s/clip) to the
un-logged clip counts. The VLM plausibility-scoring pass ($630$ clips
scored in one job) is estimated separately at $\sim\!0.3$\,h from that
job's own log-file timestamps (start to last output write, since it did
not log \texttt{elapsed\_seconds} either). This total does not include
retry/gap-fill overhead from failed job attempts (several jobs across the
project failed on environment or timeout bugs before a working
configuration was found, and were re-run) -- those wasted GPU-hours are
not reconstructable from the final result files, which only reflect the
successful run.

% ---------------------------------------------------------------------------
\section{Reproducibility checklist}
\label{app:checklist}
% ---------------------------------------------------------------------------
WACV's author guide points to the Reproducibility Checklist as a guide rather
than a mandated form, and encourages voluntary code submission. We record the
substantive items here.

\begin{itemize}[leftmargin=1.2em,itemsep=1pt,topsep=2pt]
\item \textbf{Code.} Metrics, bootstrap CIs, mechanism selection, and the
synthetic smoke test (\texttt{physweep\_metrics.py}), the system
generators (\texttt{physweep/render.py}), the tracker
(\texttt{physweep/track.py}, both blob-centroid and the CoTracker3
wrapper), and every experiment driver script (\texttt{experiments/*.py},
one per experiment plus the consolidated plausibility-and-robustness
harness) exist in the project repository, are anonymizable, and will be
released in full upon acceptance.
\item \textbf{Data.} Fully procedural, no external assets, no human
annotation, seeded and deterministic.
\item \textbf{Models.} All public, all frozen, inference only. No training or
fine-tuning was performed at any point.
\item \textbf{Randomness.} Every generation is keyed to an explicit integer
seed; $m\ge5$ seeds per grid point; seeds are recorded per clip.
\item \textbf{Statistics.} All intervals are $95\%$ paired bootstraps over
$(\theta_i,\text{seed})$; effect sizes are reported alongside $p$-values;
multiplicity across the six paired-bootstrap extrapolation tests is handled
by a Holm correction.
\item \textbf{Raw outputs.} Per-clip records (model, system, seed, $\theta$,
$\thetah$, fit $R^2$, in/out-of-range split, and, for the plausibility
pass, the VLM score) exist as JSON files in the project's \texttt{results/}
directory, one file per model-system pair plus the plausibility-scored
variants -- every number in this paper and its supplementary tables was
recomputed directly from these files during finalization, not carried
forward from memory. These files will be released in full upon
acceptance, consolidated into a single per-clip CSV.
\end{itemize}

% ---------------------------------------------------------------------------
\section{Extended related work}
\label{app:related}
% ---------------------------------------------------------------------------
This section extends the main paper's Related Work with material not
central enough to the paper's own contribution to justify main-text space,
but relevant to situating it.

\paragraph{Mechanism and memorization.} \citet{kang2025howfar} find
in-distribution success with out-of-distribution case-based behavior
(the source of the case-clamping hypothesis we test), and the
memorization literature links such collapse to overestimated training
modes amplified by classifier-free guidance \citep{diffmemorize2025}; we
turn these into a testable dichotomy (the main paper's Problem section)
and a guidance-scale ablation (the main paper's Results section).

\paragraph{Adjacent but distinct measurement approaches.}
Property-readout work recovers physical quantities from \emph{input}
videos via internal features \citep{inferring2025dynamic} or probes
diffusion states \citep{invisiblehand2026}, characterizing the
\emph{encoder} where we measure generation without weights or
activations. VLMs are separately used to plan plausible generation
\citep{vlipp2025} and judge implausibility \citep{travl2025}, and recent
analyses question whether they reason about physical transformations at
all \citep{vlmcannotreason2026} -- consistent with our own finding (the
main paper's Results section) that a VLM plausibility rater does not
track physical coherence on our stimuli. Classic benchmarks
\citep{bear2021physion,bakhtin2019phyre} evaluate prediction or planning
rather than generation, and quality metrics
\citep{unterthiner2018fvd,huang2024vbench} are orthogonal to all of the
above: they score visual fidelity, not parameter faithfulness.

{
    \small
    \bibliographystyle{ieeenat_fullname}
    \bibliography{main}
}

\end{document}
